\documentclass[10pt,a4paper]{article}
\usepackage[utf8]{inputenc}
\usepackage[T1]{fontenc}
\usepackage[french]{babel}
\usepackage{lmodern}
\usepackage[margin=1.0cm]{geometry}
\usepackage{xcolor}
\usepackage{amsmath,amssymb}
\usepackage{enumitem}
\usepackage[most]{tcolorbox}
\usepackage{hyperref}
\hypersetup{colorlinks=false,pdfborder={0 0 0}}
\setlength{\parindent}{0pt}
\setlength{\parskip}{6pt}
\renewcommand{\baselinestretch}{1.04}
\setlist[itemize]{leftmargin=*,topsep=3pt,itemsep=4pt,parsep=2pt}
\setlist[enumerate]{leftmargin=*,topsep=3pt,itemsep=4pt,parsep=2pt}
\tcbset{enhanced,boxrule=.45pt,arc=1.2mm,left=3.2mm,right=3.2mm,top=3.6mm,bottom=3.6mm,boxsep=.9mm,breakable,
  colback=white,colframe=black!48,colbacktitle=white,coltitle=black,fonttitle=\bfseries\small,
  before skip=6.5pt,after skip=6.5pt,before upper={\setlength{\parskip}{6pt}},boxed title style={boxrule=0pt,colframe=white,colback=white,left=0pt,right=0pt,top=0pt,bottom=1pt}}
\newtcolorbox{boite}[1]{title={#1}}
\newcommand{\ligne}{\vspace{3.5pt}}
\begin{document}
\thispagestyle{empty}

{\large\bfseries Feuille-méthode --- Utiliser le point fixe pour déterminer une limite}\hfill {\small Terminale spécialité}

\begin{boite}{Objectif}
Pour une suite récurrente de la forme \(u_{n+1}=f(u_n)\), la méthode du point fixe sert à \textbf{identifier la limite possible}.

\textbf{Attention :} elle s'utilise seulement \textbf{après avoir établi la convergence} de la suite.
\end{boite}

\noindent\begin{minipage}[t]{0.48\textwidth}
\begin{boite}{Méthode en étapes}
On suppose que \(u_{n+1}=f(u_n)\).

\ligne
\begin{enumerate}
  \item Montrer ou rappeler que \((u_n)\) est convergente.\newline
        On note \(\ell\) sa limite.

  \item Vérifier que \(f\) est continue sur l'intervalle où se trouvent les termes de la suite.

  \item Passer à la limite dans \(u_{n+1}=f(u_n)\).

  \item Obtenir l'équation de point fixe :
  \[
  \ell=f(\ell).
  \]

  \item Résoudre cette équation, puis garder la solution compatible avec le contexte.
\end{enumerate}
\end{boite}

\begin{boite}{Pièges fréquents}
\begin{itemize}
  \item Le point fixe \textbf{ne prouve pas} la convergence.

  \item Il faut nommer la limite :\newline
  « On note \(\ell\) la limite de \((u_n)\). »

  \item Il faut citer la \textbf{continuité} de \(f\).

  \item S'il y a plusieurs solutions, on utilise l'intervalle des termes pour conclure.
\end{itemize}
\end{boite}
\end{minipage}
\hfill
\begin{minipage}[t]{0.49\textwidth}
\begin{boite}{Rédaction type}
La suite \((u_n)\) est convergente ; on note \(\ell\) sa limite.

La fonction \(f\) est continue sur l'intervalle considéré.

Comme \(u_n\to\ell\), on a \(f(u_n)\to f(\ell)\).

\ligne
Or \(u_{n+1}=f(u_n)\), et \((u_{n+1})\) a la même limite que \((u_n)\). Donc :
\[
\ell=f(\ell).
\]

On résout cette équation et on conserve la solution compatible avec le contexte.
\end{boite}

\begin{boite}{Bilan à retenir}
Phrase clé :

\begin{center}
\fbox{\begin{minipage}{0.88\linewidth}
Par continuité de \(f\), on passe à la limite dans \(u_{n+1}=f(u_n)\).
\end{minipage}}
\end{center}

La méthode donne une limite \textbf{possible}, une fois la convergence établie.
\end{boite}
\end{minipage}

\begin{boite}{Exemple corrigé complet}
On considère la suite \((u_n)\) définie par \(u_0=1\) et \(u_{n+1}=\sqrt{2+u_n}\). On admet que \((u_n)\) est convergente et que, pour tout \(n\), \(1\leq u_n\leq 2\). Déterminer sa limite.

\textbf{Correction.}

\ligne
La suite \((u_n)\) est convergente ; on note \(\ell\) sa limite.

La fonction \(f:x\mapsto\sqrt{2+x}\) est continue sur \([1,2]\).

\ligne
Comme \(u_n\to\ell\), on a \(f(u_n)\to f(\ell)\).

\ligne
Or \(u_{n+1}=f(u_n)\), et \((u_{n+1})\) a la même limite que \((u_n)\). Donc :
\[
\ell=f(\ell),\qquad \text{c'est-à-dire}\qquad \ell=\sqrt{2+\ell}.
\]

\ligne
Comme \(\ell\in[1,2]\), on résout : \(\ell^2=2+\ell\), donc \((\ell-2)(\ell+1)=0\). On conserve la solution compatible avec \([1,2]\) : \(\boxed{\ell=2}\).
\end{boite}

\begin{boite}{Exemple guidé}
Soit \((u_n)\) définie par \(u_0=0\) et \(u_{n+1}=\dfrac12u_n+1\). On a déjà montré que \((u_n)\) est croissante et majorée par \(2\).

Compléter la rédaction en suivant exactement le modèle : convergence, notation de la limite, continuité, passage à la limite, équation de point fixe, résolution, conclusion.
\end{boite}

\end{document}
